% Chapter 4

\chapter{روش‌شناسی و معماری پیشنهادی}

\section{رویکرد کلی}
طراحی (conceptual)، پیاده‌سازی (prototype)، ارزیابی (experimental با معیارهای RMSE برای RUL).

\section{معماری پیشنهادی دوقلوی دیجیتال چندعاملی}
جزئیات ۵ لایه (موجودیت، جمع‌آوری، ادغام، DT، کاربرد) با دیاگرام‌های UML، flowchart، و شکل‌های تعاملی.
 (جزئیات لایه‌ها به همراه شکل و دیاگرام جهت واضح‌تر شدن)


\section{توصیف مجموعه داده \lr{CMAPSS}}

ساختار (۴ زیرمجموعه FD001-FD004)، حجم، پیش‌پردازش (normalization, feature selection)، و چالش‌ها (noise در داده‌ها).

\section{بررسی معماری پیشنهادی برای مجموعه داده \lr{CMAPSS}}

(بررسی عامل‌های لایه کاربرد و کارکرد آن‌ها و توضیح انتراعی آن‌ها و ارتباط آن‌ها - ارتباط با مدل‌های زبانی - پایش بی‌درنگ - پیش‌بینی و ویژگی‌های آن - ...)


عامل‌های لایه کاربرد: توضیح abstract agents (مانند monitoring agent, prediction agent).
ارتباط با LLM: ادغام برای پرس‌وجو، کنترل (خاموش/روشن)، و توضیح دلایل (با مثال promptها در Dify).
پایش بی‌درنگ: شبیه‌سازی stream data.
پیش‌بینی: ویژگی‌های LSTM (hyperparameters, training).


\section{جمع‌بندی}
مزایای معماری در مقایسه با کارهای قبلی.



%
%\chapter{فرمول}
%\label{chap:propml}
%
%نمونه‌های فرمول چندخطی:
%\begin{align}
%Loss(t,y(x)) & = \{y(x)-t\}^2 \nonumber\\ 
%& = \{y(x)- \mathbb{E}_{t}[t|x] + \mathbb{E}_{t}[t|x] -t\}^2 \nonumber \\ 
%& = \{y(x)- \mathbb{E}_{t}[t|x]\}^2 + 2\{y(x)- \mathbb{E}_{t}[t|x]\}\{\mathbb{E}_{t}[t|x] - t\} + 
%\{\mathbb{E}_{t}[t|x] -t\}^2 
%\end{align}
%برای رفرنس دهی به رابطه \eqref{eq:biasvar_bias_var} اینگونه عمل کنید.
%\begin{multline}
%\label{eq:biasvar_bias_var}
%\mathbb{E}_\mathcal{D} [\{y(x;\mathcal{D})-\mathbb{E}_{t}[t|x]\}^2] \\
%=\{\mathbb{E}_\mathcal{D}[y(x;\mathcal{D})]-\mathbb{E}_{t}[t|x]\}^2 + 
%\mathbb{E}_\mathcal{D}[\{y(x;\mathcal{D})-\mathbb{E}_\mathcal{D}[y(x;\mathcal{D})]\}^2] 
%\end{multline}
%برای نوشتن فرمول چندبخشی اینگونه عمل کنید.
%\begin{subequations}
%	\begin{equation}
%	\text{\gls{variance}} = \int{\mathbb{E}_\mathcal{D}[\{y(x;\mathcal{D})-\mathbb{E}_\mathcal{D}[y(x;\mathcal{D})]\}^2]p(x)dx}
%	\end{equation}
%	\begin{equation}
%	\text{\gls{bias}} = \int{\{\mathbb{E}_\mathcal{D}[y(x;\mathcal{D})]-\mathbb{E}_{t}[t|x]\}^2p(x)dx}
%	\end{equation}
%	\begin{equation}
%	\text{\gls{noise}} = \int{\{\mathbb{E}_{t}[t|x] -t\}^2p(x)dx}
%	\end{equation}
%\end{subequations}
%دیگر نمونه فرمول ها:
%%\begin{equation*}
%G=
%%\begin{bmatrix}
%%G_{11}&G_{12}\\G_{21}&G_{22}
%%\end{bmatrix} 
%%\end{equation*}
%
%%\begin{equation*}
%%\left\{\begin{array}{l}
%%z=G_{11}w+G_{12}u\\
%%y=G_{21}w+G_{22}u
%%\end{array}\right. 
%%\end{equation*}
%
%برای رسم دیاگرام به صورت مستقیم، اینگونه عمل کنید.
%\setlength{\unitlength}{1cm}
%\begin{figure}[t]
%	\centering
%	\lr{
%		\begin{picture}(10.5,2.3)(0,0)
%		\put(0,1.5){\vector(1,0){1.35}}
%		\put(0.675,1.7){\makebox(0,0)[b]{$r(t)$}}
%		\put(1.2,1.7){\makebox(0,0){$ \scriptstyle + $}}
%		\put(1.5,1.5){\circle{0.3}}
%		\put(1.65,1.5){\vector(1,0){1.35}}
%		\put(2.325,1.7){\makebox(0,0)[b]{$e(t)$}}
%		\put(3,1){\framebox(1.5,1){$K(t) $}}
%		\put(4.5,1.5){\vector(1,0){1.5}}
%		\put(5.25,1.7){\makebox(0,0)[b]{$u(t)$ }}
%		\put(6,1){\framebox(1.5,1){$G(s)$}}
%		\put(7.5,1.5){\line(1,0){1.5}}
%		\put(9,1.5){\circle*{0.08}}
%		\put(9,1.5){\vector(1,0){1.5}}
%		\put(9.75,1.7){\makebox(0,0)[b]{$y(t)$}}
%		\put(9,1.5){\line(0,-1){1.5}}
%		\put(9,0){\line(-1,0){7.5}}
%		\put(1.5,0){\vector(0,1){1.35}}
%		\put(1.3,1.2){\makebox(0,0){$ \scriptstyle - $}}
%		\end{picture}
%	}
%	\caption{یک سیستم زمان‌پیوسته}
%	\label{fig:sample}
%\end{figure} 
%
%
